\chapter{Conclusion}
\label{sec:conclusion}

This thesis implemented and evaluated four regular-expression parsers
over a shared core language and frontend. Two follow Brzozowski
derivatives and two follow Antimirov partial derivatives
(Chapters~\ref{sec:derivative-parser} and
\ref{sec:partial-derivative-parser}), each family with a two-pass
injection variant and a bit-coded variant. All four operate on the
same six-constructor \texttt{Regex} from Chapter~\ref{sec:frontend},
measured on the two datasets of Chapter~\ref{sec:evaluation}.

The measurements separate the four parsers into two groups.
Partial-derivative parsers are fastest and lightest on the real
corpus. Derivative parsers are slowest and heaviest, since both retain
every intermediate expression until the backward pass, so their live
state grows with input length.

Within the derivative family, bit-coding wins unconditionally.
\texttt{deriv\_bc} carries bounded state where both
\texttt{deriv\_std} variants carry the whole derivation sequence, so
it is faster and lighter at every scale measured. On
multi-million-character \texttt{(a|b|ab)*} inputs at the scale of
\citet{sulzmannlu2014}'s own benchmark, both standard variants fail
outright at every input length tested, by two independent mechanisms,
stack overflow for the recursive variant and heap exhaustion for the
iterative one, while \texttt{deriv\_bc} completes the full
ten-million-character range in under a minute.

Within the partial-derivative family the same claim does not hold.
\texttt{pderiv\_std} and \texttt{pderiv\_bc} trade leading position
with input structure rather than size. \texttt{pderiv\_std} is
marginally faster on the real corpus and on \texttt{(a|b|ab)*} over
\texttt{a}-only inputs, while \texttt{pderiv\_bc} is consistently
faster on the same pattern over alternating \texttt{ab} inputs. The
choice between them is therefore per-pattern and per-input.

Construction therefore matters more than variant. A POSIX parser
that scales to long inputs has one answer, \texttt{deriv\_bc}, the
only POSIX variant with bounded state. A non-POSIX parser must be
chosen by measurement. Where the input is short and a full POSIX tree
is not needed, \texttt{deriv\_std\_loop} remains usable on
code-complexity grounds, and \texttt{deriv\_std\_rec} stays closest
to the published algorithm.

All five parsers correctly classify every case in both datasets.
POSIX and non-POSIX tree shapes diverge only on ambiguous patterns and
reflect a policy difference, not an implementation error.

Three directions remain open. The bit-coded representation is
compatible with incremental parsing, since the bit vector that
reconstructs a tree can equally well serve as the state from which a
subsequent input character is consumed. The frontend's uniform 1000
bound cap (Section~\ref{sec:frontend-bounds}) could be replaced with a
per-pattern cost analysis. The remaining multi-million-character
sweeps in Section~\ref{sec:eval-timing-pathological} await their
configured scale.